% The four parts of the Day 8 lab. The steps follow Labs/day08/README.md.
%
% Hardware status: not tested on hardware. A value that only the board can
% give is written as "..." and the students measure it. The deck does not
% show a result that the students must predict.
%
% The frame "Part B: how to read the comparison" gives the answer of the
% second prediction of Part B. Show it after section 7 of the notebook.

% ------------------------------------------------------------------- Part A
\begin{frame}[fragile]{Part A: pre-trained models (30 min)}
  \small
  \renewcommand{\arraystretch}{1.15}
  \begingroup\centering
  \begin{tabular}{@{}L{0.07\textwidth}L{0.70\textwidth}L{0.12\textwidth}@{}}
    \toprule
    \textbf{Step} & \textbf{Work} & \textbf{Time} \\
    \midrule
    1 & Start the board, copy two folders, and connect & 8 min \\
    2 & Get two test images: \code{bus.jpg}, and a photo of your desk
      & 4 min \\
    3 & The SSD model, with Task A1 & 9 min \\
    4 & The YOLO model: two image sizes, two IoU thresholds & 9 min \\
    \bottomrule
  \end{tabular}\par\endgroup
\begin{lstlisting}[style=shell]
ssh edge@pi-NN.local "mkdir -p ~/edgeai/day08/images"    # on the laptop
scp -r day08/pi day08/models edge@pi-NN.local:~/edgeai/day08/
ssh edge@pi-NN.local                                     # then on the board
cd ~/edgeai/day08
source ~/yolo/bin/activate
rpicam-jpeg --output images/desk.jpg --width 1280 --height 960
\end{lstlisting}
  \begin{itemize}
    \item Run the first two commands in the folder \code{Labs/}. \code{NN}
          is the number of your group.
    \item \code{bus.jpg} is the image of the lecture. Step 2 of the README
          gives the command that copies it from the package.
    \item Put the bottle and the cup in front of the camera for the photo.
  \end{itemize}
\end{frame}

\begin{frame}[fragile]{Part A: the SSD model and Task A1}
  \small
  \code{python pi/detect\_ssd.py images/bus.jpg} runs SSD MobileNet V1 of
  the kit lab:
\begin{lstlisting}
Input:  shape [1, 300, 300, 3], type uint8
Output 0: shape [1, 10, 4], type float32
First inference: ... ms. Median of the next 20 runs: ... ms.
Rows in the output: 10. Rows with a score of 0.50 or more: ...
Task A1: not complete
\end{lstlisting}
  \begin{columns}[c]
    \begin{column}{0.50\textwidth}
      \begin{itemize}
        \item \textbf{Task A1:} the function \code{to\_pixel\_boxes}. A
              box of the model is \code{ymin, xmin, ymax, xmax} from 0 to
              1. Change it to pixels of the image. Keep the boxes above
              the score threshold.
        \item The diagram is the check of the script.
        \item \textbf{You record:} the four outputs, the times, and the
              detections.
      \end{itemize}
    \end{column}
    \begin{column}{0.46\textwidth}
      \centering
      \begin{tikzpicture}[font=\scriptsize, x=1cm, y=-1cm,
          >={Stealth[length=4pt]}]
        % The output of the model: a square from 0 to 1.
        \draw[coolgray, fill=boxgray] (0,0) rectangle (1.6,1.6);
        \draw[cardinalred, thick] (0.16,0.40) rectangle (0.80,1.20);
        \node[coolgray, anchor=south] at (0.8,0) {0 to 1};
        \node[anchor=north, align=center] at (0.8,1.65)
          {y: 0.25 to 0.75\\x: 0.10 to 0.50};
        \draw[->] (1.75,0.8) -- (2.35,0.8);
        % The image: 200 x 100 pixels.
        \draw[coolgray, fill=boxgray] (2.5,0.2) rectangle (4.9,1.4);
        \draw[cardinalred, thick] (2.74,0.5) rectangle (3.7,1.1);
        \node[coolgray, anchor=south] at (3.7,0.2) {200 $\times$ 100 pixels};
        \node[anchor=north, align=center] at (3.7,1.65)
          {y: 25 to 75\\x: 20 to 100};
      \end{tikzpicture}
    \end{column}
  \end{columns}
\end{frame}

\begin{frame}[fragile]{Part A: the YOLO model}
  \small
\begin{lstlisting}[style=shell]
python pi/detect_yolo.py images/bus.jpg               # YOLO11n, 640 pixels
python pi/detect_yolo.py images/bus.jpg --imgsz 320
python pi/detect_yolo.py images/bus.jpg --iou 0.9
python pi/detect_yolo.py images/desk.jpg --output yolo_desk.jpg
\end{lstlisting}
\begin{lstlisting}
Thresholds: score 0.25, IoU 0.70. Image size: 640.
Load: ... ms. First run: ... ms.
Median of the next 5 runs, in ms:
  pre-processing ...   inference ...   post-processing ...   complete call ...
Detections: ...
\end{lstlisting}
  \begin{exerciseblock}
    \textbf{Prediction,} before the second and the third command:
    \begin{itemize}
      \item The image size goes from 640 to 320 pixels. By which factor
            does the inference time change?
      \item The IoU threshold goes from 0.7 to 0.9. Do you get more
            detections, or fewer?
    \end{itemize}
  \end{exerciseblock}
  \begin{itemize}
    \item The first run downloads the weights \code{yolo11n.pt} (5.4 MB).
    \item \textbf{You record:} the time of each step, the detections, and
          the result for the photo of your desk with each model.
  \end{itemize}
\end{frame}

% ------------------------------------------------------------------- Part B
\begin{frame}{Part B: custom model (50 min)}
  \small
  Work on the laptop, in the notebook \code{custom\_detector.ipynb}. The
  bar shows its sections.
  \vskip 0.3em
  \begingroup\centering
  \begin{tikzpicture}[font=\scriptsize, x=0.2cm, y=1cm]
    % x unit: 1 = one minute. The bar is 50 minutes long.
    \fill[coolgray!35]   (0,0)  rectangle (8,0.6);
    \fill[treegreen!60]  (8,0)  rectangle (10,0.6);
    \fill[darkcyan!45]   (10,0) rectangle (17,0.6);
    \fill[darkcyan!45]   (17,0) rectangle (25,0.6);
    \fill[darkcyan!45]   (25,0) rectangle (35,0.6);
    \fill[coolgray!35]   (35,0) rectangle (43,0.6);
    \fill[coolgray!35]   (43,0) rectangle (50,0.6);
    \foreach \t in {8,10,17,25,35,43} {
      \draw[white, line width=1pt] (\t,0) -- (\t,0.6);
    }
    \node at (4,0.3)    {0, 1: data};
    \node at (9,0.3)    {2};
    \node at (13.5,0.3) {3: Task 1};
    \node at (21,0.3)   {4: Task 2};
    \node at (30,0.3)   {5: Task 3};
    \node at (39,0.3)   {6: curves};
    \node at (46.5,0.3) {7: test};
    \foreach \t in {0,8,10,17,25,35,43,50} {
      \draw[coolgray] (\t,0) -- (\t,-0.1);
    }
    \foreach \t in {0,10,17,25,35,43} {
      \node[below, text=coolgray] at (\t,-0.08) {\t};
    }
    \node[below, text=coolgray] at (50,-0.08) {50 min};
    % The training runs in the background.
    \fill[treegreen!60] (10,-0.95) rectangle (19,-0.6);
    \node[anchor=west] at (19.5,-0.775)
      {the training in the background: about 9 minutes};
  \end{tikzpicture}\par\endgroup
  \vskip 0.4em
  \renewcommand{\arraystretch}{1.1}
  \begin{tabular}{@{}lrrrL{0.50\textwidth}@{}}
    \toprule
    \textbf{Split} & \textbf{Images} & \code{bottle} & \code{cup}
      & \textbf{Use} \\
    \midrule
    \code{train} & 240 & 172 & 185 & The training changes the weights. \\
    \code{valid} & 60  & 41  & 43  & Selects the best epoch. Gives the
                                     calibration images for \code{int8}. \\
    \code{test}  & 100 & 71  & 77  & The numbers of your report. \\
    \bottomrule
  \end{tabular}
  \vskip 0.3em
  \begin{itemize}
    \item Section 1 downloads the 400 images (about 60 MB). Section 2
          starts the training. Continue with Task 1 while it runs.
    \item The cell after a task prints \code{Task N: complete} or
          \code{not complete}.
    \item \textbf{You record} (section 1): does each box touch its object
          on the four sides? Which object has no label?
  \end{itemize}
  \source{Images and labels: COCO 2017, \texttt{cocodataset.org}.}
\end{frame}

\begin{frame}{Part B: Tasks 1, 2, and 3}
  \small
  \renewcommand{\arraystretch}{1.2}
  \begingroup\centering
  \begin{tabular}{@{}L{0.07\textwidth}L{0.19\textwidth}L{0.64\textwidth}@{}}
    \toprule
    \textbf{Task} & \textbf{Function} & \textbf{You write} \\
    \midrule
    1 & \code{box\_iou} & The IoU of two boxes: the intersection, the
        union, and their ratio \\
    2 & \code{nms} & The suppression for one class: keep the box with the
        highest score, remove each box with an IoU above the threshold,
        and do this again \\
    3 & \code{count\_image} & TP, FP, and FN for one class in one image:
        each detection takes the free label with the highest IoU \\
    \bottomrule
  \end{tabular}\par\endgroup
  \vskip 0.3em
  \begin{columns}[c]
    \begin{column}{0.30\textwidth}
      \centering
      \begin{tikzpicture}[font=\scriptsize, x=0.17cm, y=-0.17cm]
        \fill[darkorange!35] (2,0) rectangle (10,10);
        \draw[treegreen, dashed, thick] (0,0) rectangle (10,10);
        \draw[cardinalred, thick] (2,0.15) rectangle (12,10.15);
        \node at (6,5) {80};
        \node[anchor=north, align=center] at (6,10.6)
          {IoU $= 80 / 120$\\$= 0.67$};
      \end{tikzpicture}
    \end{column}
    \begin{column}{0.66\textwidth}
      The checks:
      \begin{itemize}
        \item Task 1: the two box pairs of the lecture, and 200 random
              pairs. The diagram shows pair 1.
        \item Task 2: \code{bus.jpg} has 47 candidates above the score
              0.25, and 5 detections for the IoU threshold 0.7, as in the
              lecture.
        \item Task 3: five small cases.
      \end{itemize}
    \end{column}
  \end{columns}
  \vskip 0.5em
  \textbf{You record:} the IoU of three pairs, and the number of
  detections for the two IoU thresholds.
\end{frame}

\begin{frame}{Part B: the baseline, the training, and the test}
  \small
  \renewcommand{\arraystretch}{1.2}
  \begingroup\centering
  \begin{tabular}{@{}L{0.11\textwidth}L{0.80\textwidth}@{}}
    \toprule
    \textbf{Section} & \textbf{Work} \\
    \midrule
    5 & The baseline: the pre-trained YOLO11n on the 100 test images, at
        320 and at 640 pixels. Your function counts TP, FP, and FN. \\
    6 & The result of the training: one line for each epoch, then the
        curves of the two losses and of the mAP \\
    7 & Your model on the 100 test images, and the comparison with the
        baseline. The validation images selected the best epoch, so they
        are not a fair test. \\
    \bottomrule
  \end{tabular}\par\endgroup
  \vskip 0.3em
  \begin{exerciseblock}
    \textbf{Two predictions.} Write each one in the report before you run
    the cell:
    \begin{itemize}
      \item Section 5: is the recall at 320 pixels lower than at 640
            pixels?
      \item Section 7: your model has 2 classes and saw 240 images. The
            pre-trained model has 80 classes and saw 118 000 images. Which
            model has the higher recall for the two classes?
    \end{itemize}
  \end{exerciseblock}
  \begin{itemize}
    \item \textbf{You record:} the best epoch with its mAP, the tables of
          the three models, and one false positive and one false negative
          of the figure.
  \end{itemize}
\end{frame}

\begin{frame}{Part B: how to read the comparison}
  \begingroup\centering
  \begin{tikzpicture}[font=\footnotesize,
      m/.style={draw=coolgray, rounded corners=3pt, fill=boxgray,
                minimum width=4.4cm, minimum height=1.5cm, align=center}]
    \node[m, draw=darkcyan]  (a) at (0,0)
      {\textbf{The baseline}\\pre-trained YOLO11n\\80 classes, 118 000
       images};
    \node[m, draw=treegreen] (b) at (5.2,0)
      {\textbf{Your model}\\the same network, a new head\\2 classes, 240
       images};
  \end{tikzpicture}\par\endgroup
  \vskip 0.5em
  \small
  \begin{itemize}
    \item On the work computer of the course, the baseline had the higher
          recall for the two classes. This is a normal result: 240 images
          do not pass a model that learned the two classes from 118 000
          images.
    \item A custom model is the correct tool in two cases:
      \begin{itemize}
        \item The pre-trained model does not know your classes. Example:
              the box and the wheel of the kit lab.
        \item Your images are different from the images of COCO: a
              different camera position, or a different light.
      \end{itemize}
    \item With more images of your own camera, the custom model becomes
          better than the baseline.
  \end{itemize}
  \keypoint{Measure a baseline before you train. Then you know what the
    training gave.}
\end{frame}

% ------------------------------------------------------------------- Part C
\begin{frame}[fragile]{Part C: export and deploy (40 min)}
  \small
  \renewcommand{\arraystretch}{1.15}
  \begingroup\centering
  \begin{tabular}{@{}L{0.07\textwidth}L{0.70\textwidth}L{0.12\textwidth}@{}}
    \toprule
    \textbf{Step} & \textbf{Work} & \textbf{Time} \\
    \midrule
    1 & Export: sections 8 to 10 of the notebook, on the laptop & 15 min \\
    2 & Copy the six files to the board & 3 min \\
    3 & Check your model with the photo of your desk & 5 min \\
    4 & Live detection with the camera & 12 min \\
    5 & Demonstration for the instructor & 5 min \\
    \bottomrule
  \end{tabular}\par\endgroup
\begin{lstlisting}[style=shell]
python pi/detect_yolo.py images/desk.jpg \
    --model models/cupbottle_320_ncnn_model --imgsz 320
python pi/live_detect.py --model models/cupbottle_320_ncnn_model
python pi/live_detect.py --model models/cupbottle_320_int8.tflite
python pi/live_detect.py --model models/cupbottle_320_ncnn_model --conf 0.5
\end{lstlisting}
  \begin{itemize}
    \item Step 2: use the \code{scp} command of section 10 of the notebook.
    \item Step 3: give the image size of the file with \code{--imgsz}.
          Compare the result with the pre-trained model of Part A.
    \item Step 4: stop the application with Ctrl+C before you start it
          with a different file.
  \end{itemize}
\end{frame}

\begin{frame}{Part C: six files from one model}
  \begingroup\centering
  \begin{tikzpicture}[font=\footnotesize, >={Stealth[length=5pt]},
      f/.style={draw=coolgray, rounded corners=2pt, fill=boxgray,
                minimum width=2.0cm, minimum height=0.62cm, align=center},
      q/.style={f, draw=cardinalred, fill=cardinalred!10},
      lab/.style={font=\scriptsize, text=coolgray, align=center}]
    \node[f, draw=treegreen, minimum height=1.0cm] (m) at (0,-0.45)
      {your model\\(PyTorch)};
    \node[lab] at (2.25,0) {320 pixels};
    \node[lab] at (2.25,-0.9) {640 pixels};
    \foreach \y in {0,-0.9} {
      \node[f] at (4.2,\y) {NCNN};
      \node[f] at (6.4,\y) {LiteRT};
      \node[q] at (8.6,\y) {LiteRT};
    }
    \draw[->] (m.east) -- (1.5,0);
    \draw[->] (m.east) -- (1.5,-0.9);
    \node[lab] at (4.2,0.65) {\texttt{float32}};
    \node[lab] at (6.4,0.65) {\texttt{float32}};
    \node[lab] at (8.6,0.65) {\texttt{int8}};
  \end{tikzpicture}\par\endgroup
  \vskip 0.3em
  \small
  \begin{itemize}
    \item A name gives the three properties:
          \code{cupbottle\_320\_int8.tflite}.
    \item Section 9 measures the mAP of each LiteRT file. An NCNN folder
          has the same \code{float32} weights: use the value of the LiteRT
          \code{float32} file.
    \item The test set has 100 images. A difference of 0.02 of the mAP is
          not a clear result.
  \end{itemize}
  \begin{exerciseblock}
    \textbf{Three predictions,} before the cells of sections 8 and 9:
    \begin{itemize}
      \item How large is the \code{int8} file, as a part of the
            \code{float32} file?
      \item How many points of mAP50-95 does the \code{int8} file lose?
      \item Your model learned at 320 pixels. Is its mAP at 640 pixels
            higher, the same, or lower?
    \end{itemize}
  \end{exerciseblock}
\end{frame}

\begin{frame}[fragile]{Part C: the live detection}
  \begingroup\centering
  \begin{tikzpicture}[font=\scriptsize, >={Stealth[length=4pt]},
      s/.style={draw=coolgray, rounded corners=2pt, fill=boxgray,
                minimum width=1.3cm, minimum height=0.6cm, align=center}]
    \node[s, draw=darkcyan] (cam) at (0,0) {camera};
    \node[s] (c) at (1.75,0) {capture};
    \node[s] (p) at (3.3,0)  {pre};
    \node[s, draw=cardinalred] (i) at (4.85,0) {inference};
    \node[s] (o) at (6.4,0)  {post};
    \node[s] (d) at (7.95,0) {draw};
    \node[s, draw=darkorange] (w) at (7.95,-1.5) {web server};
    \node[s, draw=darkcyan] (b) at (3.3,-1.5) {browser of the laptop};
    \foreach \a/\b in {cam/c, c/p, p/i, i/o, o/d} { \draw[->] (\a) -- (\b); }
    \draw[->] (d) -- node[right, coolgray, pos=0.72] {JPEG} (w);
    \draw[->] (w) -- node[above, coolgray] {port 5000, Wi-Fi} (b);
    \draw[coolgray, dashed] (0.95,0.5) rectangle (8.95,-0.5);
    \node[coolgray] at (4.95,0.7)
      {one thread on the board: one frame after the other};
  \end{tikzpicture}\par\endgroup
\begin{lstlisting}[basicstyle=\ttfamily\tiny]
Model: cupbottle_320_ncnn_model, NCNN, float32, image size 320
Open http://<name of this board>:5000 in a browser. Stop with Ctrl+C.
  ... frames/s | capture   ...  pre  ...  inference    ...  post  ...  draw  ... ms | 1 bottle, 1 cup | ... C
\end{lstlisting}
  \small
  \begin{itemize}
    \item Open \code{http://pi-NN.local:5000} on the laptop. You see the
          live image with the boxes, the frame rate, and the time of each
          step.
    \item Put the bottle and the cup in front of the camera. Move them,
          change the distance, and change the background.
    \item Then start the \code{int8} file, and the first file with the
          score threshold 0.5.
    \item \textbf{You record:} the frame rate and the inference time of
          each command, the objects that the model finds, the objects that
          it misses, and the wrong detections.
  \end{itemize}
\end{frame}

% ------------------------------------------------------------------- Part D
\begin{frame}[fragile]{Part D: measure (30 min)}
  \small
  \renewcommand{\arraystretch}{1.15}
  \begingroup\centering
  \begin{tabular}{@{}L{0.07\textwidth}L{0.70\textwidth}L{0.12\textwidth}@{}}
    \toprule
    \textbf{Step} & \textbf{Work} & \textbf{Time} \\
    \midrule
    1 & Predict: write your answer for each question of the report
      & 5 min \\
    2 & Task D1: the function \code{frame\_rate} in
        \code{pi/bench\_detect.py} & 5 min \\
    3 & Measure: run the script two times & 10 min \\
    4 & Fill the frame-rate table, and calculate the ratios & 5 min \\
    5 & Compare with your predictions, and write the Decision Log
      & 5 min \\
    \bottomrule
  \end{tabular}\par\endgroup
\begin{lstlisting}[basicstyle=\ttfamily\tiny]
model                        runtime type     size | capture     pre inference    post | frames/s | objects
cupbottle_320_ncnn_model     NCNN    float32   320 |     ...     ...       ...     ... |      ... | ...
...
cupbottle_640_int8.tflite    LiteRT  int8      640 |     ...     ...       ...     ... |      ... | ...
\end{lstlisting}
  \begin{itemize}
    \item \code{python pi/bench\_detect.py} measures the six files: the
          median time of each step for 30 frames of the camera. It also
          writes \code{results.csv}.
    \item Task D1: one frame needs the sum of the times of the four steps.
          The check is the example of the lecture: 325 ms for one frame
          give 3.1 frames in each second.
  \end{itemize}
\end{frame}

\begin{frame}{Part D: predict, measure, compare}
  \begin{exerciseblock}
    \textbf{Prediction,} before you run the script:
    \begin{itemize}
      \item \code{float32}: which runtime is faster, NCNN or LiteRT?
      \item The gain of the frame rate from 640 to 320 pixels: 4, more
            than 4, or less than 4?
      \item LiteRT: is \code{int8} faster than \code{float32}? By which
            factor?
      \item Which step is the largest after the inference?
    \end{itemize}
  \end{exerciseblock}
  \small
  \begin{itemize}
    \item Stop the live application. Point the camera at the bottle and
          the cup. Use the second run for the table: the difference
          between the two runs is the noise of your measurement.
    \item \textbf{You record:} the six rows with the mAP50-95 of section 9,
          the two temperatures, and the ratios: 320 against 640 pixels,
          \code{int8} against \code{float32}, NCNN against LiteRT.
    \item The lecture measured on an x86 processor: \code{int8} was slower
          than \code{float32}, and NCNN was slower than LiteRT. What do
          you measure?
  \end{itemize}
\end{frame}
